\documentclass[10pt,a4paper]{article}
\usepackage[utf8]{inputenc}
\usepackage[T1]{fontenc}
\usepackage[english]{babel}
\usepackage{geometry}
\geometry{margin=1.8cm}
\usepackage{enumitem}
\usepackage{booktabs}
\usepackage{fancyhdr}
\usepackage{lastpage}
\usepackage{tcolorbox}
\usepackage{hyperref}

\pagestyle{fancy}
\fancyhf{}
\rfoot{\thepage/\pageref{LastPage}}
\lhead{\small TOEFL ITP 2026 --- Day 3: Structure Deep Dive}
\rhead{\small MPrometeo}
\renewcommand{\headrulewidth}{0.4pt}
\setlength{\parindent}{0pt}
\setlength{\parskip}{5pt}

\begin{document}

\begin{center}
{\LARGE \textbf{DAY 3 --- Sunday Mar 1}}\\[4pt]
{\large Structure Deep Dive}\\[2pt]
{\normalsize Based on \emph{One Battle After Another} (2025)\\
Paul Thomas Anderson's adaptation of Pynchon's \emph{Vineland}}\\[2pt]
\rule{\textwidth}{0.8pt}
\end{center}

% ============================================================
\section*{Context: The Film \& The Novel}
% ============================================================

\begin{tcolorbox}[colback=white, colframe=black,
    title=One Battle After Another (2025)]
\textbf{Director:} Paul Thomas Anderson\\
\textbf{Based on:} \emph{Vineland} (1990) by Thomas Pynchon\\
\textbf{Cast:} Leonardo DiCaprio (Bob/``Ghetto Pat''), Benicio Del
Toro, Sean Penn (Col.\ Lockjaw), Teyana Taylor (Perfidia), Regina
Hall, Chase Infiniti (Willa)\\[4pt]
\textbf{Plot:} Bob, a stoned ex-revolutionary, has been hiding for
16 years with his daughter Willa after his partner Perfidia was
captured by the fascistic Col.\ Lockjaw. When Lockjaw resurfaces
and Willa goes missing, Bob reunites his old radical crew for one
last fight. Anderson described his process as: ``I stole the parts
that spoke to me and just started running like a thief.''\\[4pt]
\textbf{2026 Oscar Nominations (12):} Best Picture, Director, Actor
(DiCaprio), Supporting Actor (Del Toro + Penn), Supporting Actress
(Taylor), Adapted Screenplay, Cinematography, Editing, Production
Design, Casting, Original Score (Jonny Greenwood), Sound.\\[4pt]
\textbf{BAFTAs Won (6):} Best Film, Director, Adapted Screenplay,
Supporting Actor (Penn), Cinematography, Editing.\\[4pt]
\textbf{Source novel --- \emph{Vineland}:} Set in Reagan-era
Northern California, Pynchon tells the story of Zoyd Wheeler, a
hippie single father, and his daughter Prairie, caught between the
ruins of 1960s counterculture and the authoritarian Drug Enforcement
agent Brock Vond. Anderson updated the setting, changed names, and
injected his own manic energy, but kept the core: a father's love,
political paranoia, and the ghosts of American radicalism.
\end{tcolorbox}

% ============================================================
\section*{Why This Film Deserves Every Oscar}
% ============================================================

\begin{tcolorbox}[colback=white, colframe=black,
    title=A Pynchon Fan's Assessment]
\small
\textbf{Is it a good Pynchon adaptation?}

Yes and no --- and that's exactly the point.

Pynchon's novels are famously ``unadaptable.'' His prose is a
kaleidoscope of paranoia, encyclopedic digressions, and sentences
that can run for half a page. \emph{Inherent Vice} (2014), Anderson's
first Pynchon adaptation, was faithful but polarizing --- some felt
it was too confusing, others loved its hazy brilliance.

With \emph{One Battle After Another}, Anderson made a crucial
decision: he stopped trying to be faithful. Instead of translating
Pynchon page-by-page, he captured the \emph{feeling} --- the
paranoid energy, the absurd humor, the heartbreak underneath the
chaos. Anderson himself admitted: ``I stole the parts that spoke to
me and just started running like a thief.''

The result is not \emph{Vineland} on screen. It is a movie that
\emph{feels} like reading Pynchon: disorienting, funny, terrifying,
and somehow deeply moving. DiCaprio's performance as Bob --- half
burnout, half revolutionary saint --- channels the spirit of Zoyd
Wheeler without imitating him. Sean Penn's Col.\ Lockjaw is Brock
Vond reimagined as a modern American fascist. And Teyana Taylor's
Perfidia brings a depth and fury that even Pynchon's novel only
hinted at.

\textbf{Does it deserve Best Picture?} Absolutely. Not because it
perfectly adapts a novel, but because it does what the best
adaptations do: it \emph{transforms} the source material into
something that could only exist as cinema, while remaining true to
the spirit that made the original great.

The 12 nominations are earned. The BAFTAs were deserved. And if it
wins Best Picture, it will be the first Pynchon adaptation to do so
--- a sentence nobody thought would ever be written.
\end{tcolorbox}

% ============================================================
\section*{Grammar Focus: Sentence Structure}
% ============================================================

\subsection*{Rule 1: Independent vs.\ Dependent Clauses}

\begin{itemize}[nosep]
\item \textbf{Independent:} Can stand alone as a sentence.\\
``Anderson adapted Pynchon's \emph{Vineland}.''
\item \textbf{Dependent:} Cannot stand alone. Begins with a
subordinating conjunction.\\
``Although Anderson adapted \emph{Vineland}, he changed many details.''
\end{itemize}

\textbf{TOEFL trap:} A dependent clause alone = fragment = WRONG.
\begin{itemize}[nosep]
\item \textbf{Fragment:} ``Because Anderson admired Pynchon's work.''
\item \textbf{Correct:} ``Anderson made the film because he admired
Pynchon's work.''
\end{itemize}

\subsection*{Rule 2: Run-on Sentences}

Two independent clauses without proper connection:
\begin{itemize}[nosep]
\item \textbf{Run-on:} ``DiCaprio starred as Bob Penn played the
villain.''
\item \textbf{Fix 1:} ``DiCaprio starred as Bob. Penn played the
villain.''
\item \textbf{Fix 2:} ``DiCaprio starred as Bob, and Penn played the
villain.''
\item \textbf{Fix 3:} ``DiCaprio starred as Bob; Penn played the
villain.''
\end{itemize}

\subsection*{Rule 3: Relative Clauses (who/which/that/whose/where)}

\begin{center}
\begin{tabular}{ll}
\toprule
\textbf{Refers to} & \textbf{Use} \\
\midrule
People & who / whom / whose \\
Things & which / that \\
Places & where \\
Time & when \\
\bottomrule
\end{tabular}
\end{center}

``The film, \textbf{which} was based on \emph{Vineland}, received
12 Oscar nominations.''

``Anderson, \textbf{whose} previous Pynchon adaptation was
\emph{Inherent Vice}, took a freer approach this time.''

\subsection*{Rule 4: Reduced Clauses}

TOEFL loves testing the short form:
\begin{itemize}[nosep]
\item Full: ``The actor \textbf{who was nominated} for Best Actor...''
\item Reduced: ``The actor \textbf{nominated} for Best Actor...''
\item Full: ``The film \textbf{that is based} on \emph{Vineland}...''
\item Reduced: ``The film \textbf{based} on \emph{Vineland}...''
\end{itemize}

\subsection*{Rule 5: Inversion after Negative Adverbials}

When ``Not only,'' ``Rarely,'' ``Seldom,'' ``Not until,'' ``Never''
start a sentence $\to$ invert subject and auxiliary:

\begin{itemize}[nosep]
\item ``Not only \textbf{did} Anderson direct the film, but he also
wrote the screenplay.''
\item ``Rarely \textbf{has} a Pynchon adaptation received such acclaim.''
\item ``Never before \textbf{had} a film based on Pynchon won a BAFTA.''
\end{itemize}

\subsection*{Rule 6: Parallel Structure}

Elements joined by ``and,'' ``or,'' ``but'' must have the same
grammatical form:

\begin{itemize}[nosep]
\item \textbf{Wrong:} ``The film is exciting, funny, and
\textbf{has great acting}.''
\item \textbf{Right:} ``The film is exciting, funny, and
\textbf{well-acted}.''
\end{itemize}

% ============================================================
\section*{Exercises: Structure (10 questions)}
% ============================================================

\begin{enumerate}[label=\textbf{\arabic*.}, leftmargin=*]

\item Paul Thomas Anderson, \rule{2cm}{0.4pt} previous Pynchon
adaptation was \emph{Inherent Vice}, took a freer approach with
\emph{Vineland}.
\begin{enumerate}[label=(\Alph*)]
\item who \item which \item whose \item whom
\end{enumerate}

\item \rule{2cm}{0.4pt} the novel had been considered unadaptable,
Anderson managed to capture its paranoid energy on screen.
\begin{enumerate}[label=(\Alph*)]
\item Despite \item Although \item Because \item Unless
\end{enumerate}

\item The film \rule{2cm}{0.4pt} on Pynchon's \emph{Vineland}
received 12 Oscar nominations in 2026.
\begin{enumerate}[label=(\Alph*)]
\item basing \item based \item which basing \item was base
\end{enumerate}

\item Not only \rule{2cm}{0.4pt} the BAFTAs, but it is also
favored to win Best Picture at the Oscars.
\begin{enumerate}[label=(\Alph*)]
\item the film won
\item did the film win
\item the film did win
\item won the film
\end{enumerate}

\item \textit{Find the error:}
(A)\textbf{Leonardo DiCaprio}, who (B)\textbf{stars} as Bob,
(C)\textbf{gave} one of (D)\textbf{his more} memorable performances.

\item Anderson described his adaptation process as
\rule{2cm}{0.4pt} the parts he loved and running with them.
\begin{enumerate}[label=(\Alph*)]
\item steal \item stealing \item stole \item to stealing
\end{enumerate}

\item \textit{Find the error:}
Sean Penn, (A)\textbf{which} plays the villain Col.\ Lockjaw,
(B)\textbf{was nominated} for Best (C)\textbf{Supporting} Actor at
(D)\textbf{the} 2026 Oscars.

\item Neither \emph{Gravity's Rainbow} \rule{2cm}{0.4pt} ever been
successfully adapted for the screen.
\begin{enumerate}[label=(\Alph*)]
\item or \emph{Mason \& Dixon} have
\item nor \emph{Mason \& Dixon} has
\item nor \emph{Mason \& Dixon} have
\item and \emph{Mason \& Dixon} has
\end{enumerate}

\item \textit{Find the error:}
The film (A)\textbf{is} exciting, (B)\textbf{visually} stunning,
and (C)\textbf{has} a soundtrack (D)\textbf{composed} by Jonny
Greenwood.

\item Rarely \rule{2cm}{0.4pt} a literary adaptation generated so
much excitement during awards season.
\begin{enumerate}[label=(\Alph*)]
\item a film has \item has a film \item a film did \item does a film
\end{enumerate}

\end{enumerate}

\subsection*{Answer Key}

\begin{center}
\begin{tabular}{cll}
\toprule
\# & Ans & Explanation \\
\midrule
1 & C & ``Whose'' = possessive (his previous adaptation) \\
2 & B & ``Although'' + clause (has subject + verb) \\
3 & B & Reduced passive: ``(which was) based on'' \\
4 & B & Inversion: ``Not only DID the film WIN'' \\
5 & D & ``His more'' $\to$ ``his MOST'' (superlative, not comparative) \\
6 & B & ``Described as'' + gerund: ``stealing'' \\
7 & A & ``Which'' $\to$ ``WHO'' (Sean Penn is a person) \\
8 & B & Neither...nor + nearest subject singular $\to$ ``has'' \\
9 & C & Parallel structure broken: ``is exciting, stunning, and HAS'' \\
  &   & $\to$ should be ``exciting, stunning, and musically rich'' \\
10 & B & Inversion: ``Rarely HAS a film'' \\
\bottomrule
\end{tabular}
\end{center}

% ============================================================
\section*{Reading Passage ($\sim$250 words)}
% ============================================================

\begin{tcolorbox}[colback=white, colframe=black]
\small
Thomas Pynchon's novels have long been considered among the most
difficult works in American literature to adapt for the screen. His
dense, digressive prose --- filled with conspiracy theories,
encyclopedic references, and sentences that can stretch for half a
page --- seems designed to resist cinematic translation. Yet in 2025,
director Paul Thomas Anderson proved that a Pynchon adaptation could
not only work as cinema but could dominate awards season.

\emph{One Battle After Another}, loosely based on Pynchon's 1990
novel \emph{Vineland}, earned 12 Academy Award nominations and swept
six categories at the BAFTAs, including Best Film and Best Director.
Anderson's approach was deliberately unfaithful. Rather than
attempting a page-by-page translation, he extracted what he called
``the parts that spoke to me'' and rebuilt them into an original
cinematic experience. The result captures the \emph{feeling} of
Pynchon --- the paranoia, the dark humor, the political anger ---
without being enslaved to the plot.

This strategy stands in contrast to Anderson's earlier Pynchon
adaptation, \emph{Inherent Vice} (2014), which attempted closer
fidelity to the source material and divided critics and audiences
alike. The lesson, perhaps, is that the best literary adaptations
are not translations but transformations. Anderson did not put
\emph{Vineland} on screen; he made a film that could only exist
because \emph{Vineland} existed first.

For Pynchon fans, the film raises a tantalizing question: if
\emph{Vineland} --- one of his ``minor'' novels --- can produce this
level of cinema, what might a filmmaker someday do with
\emph{Gravity's Rainbow}?
\end{tcolorbox}

\subsection*{Reading Questions}

\begin{enumerate}[label=\textbf{\arabic*.}]

\item What is the main idea of the passage?
\begin{enumerate}[label=(\Alph*)]
\item Pynchon's novels are impossible to adapt
\item Anderson proved a Pynchon adaptation could succeed by
transforming rather than translating the source
\item \emph{Inherent Vice} was a better film
\item Pynchon should write screenplays
\end{enumerate}

\item ``Digressive'' in paragraph 1 is closest in meaning to:
\begin{enumerate}[label=(\Alph*)]
\item focused \item wandering from the main subject
\item boring \item simple
\end{enumerate}

\item How did Anderson's approach to \emph{One Battle After Another}
differ from \emph{Inherent Vice}?
\begin{enumerate}[label=(\Alph*)]
\item He was more faithful to the novel
\item He was deliberately less faithful, extracting only key elements
\item He didn't read the novel at all
\item He used the same cast
\end{enumerate}

\item The word ``enslaved'' in paragraph 2 most likely means:
\begin{enumerate}[label=(\Alph*)]
\item inspired by \item controlled or trapped by
\item unaware of \item confused by
\end{enumerate}

\item It can be inferred that the author believes:
\begin{enumerate}[label=(\Alph*)]
\item All Pynchon novels should be adapted exactly as written
\item Transformation is more effective than literal translation for
Pynchon adaptations
\item \emph{Vineland} is Pynchon's best novel
\item \emph{Gravity's Rainbow} will never be adapted
\end{enumerate}

\item What is ``tantalizing'' closest in meaning to?
\begin{enumerate}[label=(\Alph*)]
\item frightening \item frustrating
\item excitingly tempting \item boring
\end{enumerate}

\item How many BAFTAs did the film win?
\begin{enumerate}[label=(\Alph*)]
\item 3 \item 6 \item 10 \item 12
\end{enumerate}
\end{enumerate}

\textbf{Answers:} 1-B, 2-B, 3-B, 4-B, 5-B, 6-C, 7-B

% ============================================================
\section*{Shadowing Exercise (15 min --- Gemini Live)}
% ============================================================

\begin{tcolorbox}[colback=white, colframe=black,
    title=Tell Gemini Live:]
``Read this paragraph aloud at normal American speed. Then I repeat
sentence by sentence. Correct my pronunciation after each one.''
\end{tcolorbox}

\begin{quote}\small
``Pynchon's novels have long been considered among the most
difficult works in American literature to adapt for the screen.
Anderson's approach was deliberately unfaithful. Rather than
attempting a page-by-page translation, he extracted what he called
the parts that spoke to me and rebuilt them into an original
cinematic experience. The best literary adaptations are not
translations but transformations.''
\end{quote}

\textbf{Pronunciation targets:}
\begin{itemize}[nosep]
\item \textbf{adaptation} --- ad-ap-TAY-shun (stress on 3rd)
\item \textbf{cinematic} --- sin-eh-MAT-ik (stress on 3rd)
\item \textbf{deliberately} --- de-LIB-er-it-lee (stress on 2nd)
\item \textbf{encyclopedic} --- en-sai-clo-PEE-dik (stress on 4th)
\item \textbf{transformation} --- trans-for-MAY-shun (stress on 3rd)
\end{itemize}

% ============================================================
\section*{Bonus: Film Vocabulary for TOEFL}
% ============================================================

Words from this lesson that appear frequently in TOEFL reading:

\begin{center}
\begin{tabular}{lll}
\toprule
\textbf{Word} & \textbf{Meaning (ES)} & \textbf{TOEFL context} \\
\midrule
adaptation & adaptaci\'on & literary/scientific \\
acclaim & aclamaci\'on & critical acclaim \\
controversial & controvertido & a controversial decision \\
deliberately & deliberadamente & deliberately chosen \\
digressive & divagante & digressive prose \\
dominate & dominar & dominate the market \\
fidelity & fidelidad & fidelity to the source \\
polarizing & polarizante & a polarizing figure \\
tantalizing & tentador & tantalizing possibilities \\
transformation & transformaci\'on & undergo transformation \\
\bottomrule
\end{tabular}
\end{center}

% ============================================================
\section*{Speaking Bonus (Gemini Live --- 15 min)}
% ============================================================

\begin{tcolorbox}[colback=white, colframe=black]
Ask Gemini Live (or ChatGPT Voice):

\textit{``Let's have a conversation about movie adaptations of books.
I'll argue that One Battle After Another deserves Best Picture.
Ask me follow-up questions and correct my English.''}

\textbf{Key phrases to practice:}
\begin{itemize}[nosep]
\item ``In my opinion, the film deserves to win because...''
\item ``What makes this adaptation successful is that...''
\item ``Unlike previous attempts, Anderson chose to...''
\item ``The performances, particularly DiCaprio's, demonstrate...''
\item ``I would argue that transformation is more important than
fidelity when adapting complex literature.''
\end{itemize}
\end{tcolorbox}

\end{document}